\documentclass[conference]{IEEEtran}
\IEEEoverridecommandlockouts

\usepackage{cite}
\usepackage{amsmath,amssymb,amsfonts}
\usepackage{algorithmic}
\usepackage{graphicx}
\usepackage{textcomp}
\usepackage{xcolor}
\usepackage{booktabs}
\usepackage{hyperref}
\usepackage{url}

\def\BibTeX{{\rm B\kern-.05em{\sc i\kern-.025em b}\kern-.08em
    T\kern-.1667em\lower.7ex\hbox{E}\kern-.125emX}}

\begin{document}

\title{SimCLR on CIFAR-10: Paper Reproduction, Data Scaling Analysis, and Interactive Web Deployment\\
\thanks{Machine Learning Project --- Fall 2026 (Deliverable 3 Final Research Report). Machine Learning Paper Reproduction \& Deployment Assignment.}
}

\author{
\IEEEauthorblockN{Saalik Hussain Zaidi}
\IEEEauthorblockA{\textit{School of Computer Science} \\
\textit{Institute of Business Administration (IBA)}\\
Karachi, Pakistan \\
ERP: 36206}
\and
\IEEEauthorblockN{Mubashir Abid}
\IEEEauthorblockA{\textit{School of Computer Science} \\
\textit{Institute of Business Administration (IBA)}\\
Karachi, Pakistan \\
ERP: 36224}
\and
\IEEEauthorblockN{Muhammad Zohair Khan}
\IEEEauthorblockA{\textit{School of Computer Science} \\
\textit{Institute of Business Administration (IBA)}\\
Karachi, Pakistan \\
ERP: 33939}
}

\maketitle

\begin{abstract}
Learning robust, semantically rich visual representations without human supervision remains a foundational objective in computer vision. In this work, we conduct an end-to-end reproduction and scientific extension of SimCLR (Chen et al., ICML 2020) on the CIFAR-10 benchmark. We implement the entire contrastive framework completely from scratch in PyTorch, comprising stochastic two-view data augmentations, a modified ResNet-18 base encoder with a custom CIFAR stem omitting max-pooling, a 2-layer MLP projection head ($512 \to 512 \to 128$), and the normalized temperature-scaled cross-entropy loss (NT-Xent). Under feasibility adjustments for single-accelerator training (Apple Silicon M4 Pro GPU, mini-batch size 256, 50 pretraining epochs), our reproduced model achieves \textbf{75.23\% Top-1 Test Accuracy} and \textbf{98.71\% Top-5 Test Accuracy} under the standard frozen linear evaluation protocol on the full 10,000-image test set. Furthermore, we execute a novel data-scaling experiment evaluating the impact of unlabeled pretraining pool size ($10\% \to 25\% \to 50\% \to 100\%$) under fixed optimizer budgets, empirically demonstrating logarithmic scaling with diminishing marginal returns ($54.17\% \to 65.37\% \to 71.04\% \to 75.23\%$). Finally, we package and deploy the trained encoder and linear classifier as a responsive, Stack Overflow-themed interactive web application hosted on Hugging Face Spaces, instrumented with real-time inference latency tracking and production observability.
\end{abstract}

\begin{IEEEkeywords}
Contrastive Learning, Self-Supervised Learning, SimCLR, CIFAR-10, Visual Representation Learning, Neural Scaling Laws, Model Deployment.
\end{IEEEkeywords}

\section{Introduction and Problem Statement}
Modern deep learning has demonstrated state-of-the-art visual recognition capabilities primarily when guided by extensive human annotations. Large-scale supervised datasets such as ImageNet \cite{russakovsky2015imagenet} have driven breakthroughs in deep residual networks \cite{he2016deep}, object detection, and semantic segmentation. However, manual annotation is labor-intensive, cost-prohibitive, and often unscalable in specialized domains (e.g., medical imaging, satellite telemetry, or fine-grained industrial sorting) where domain expertise is scarce. Furthermore, human supervision forces neural networks to learn representations tied narrowly to discrete class labels, often discarding rich contextual information present in the raw visual signal.

To circumvent annotation bottlenecks, unsupervised visual representation learning aims to train general-purpose feature extractors entirely on unlabeled imagery. Once pretrained, the frozen feature extractor can transfer effectively to downstream prediction tasks using minimal labeled supervision. In this project, we reproduce and evaluate \textbf{SimCLR} (Chen et al., ICML 2020 \cite{chen2020simple}), a seminal contrastive learning framework that established that unsupervised pretraining could rival supervised pretraining on ImageNet without specialized memory banks or architectural heuristics.

\textbf{Project Scope and Objectives:}
\begin{enumerate}
    \item \textit{From-Scratch Implementation:} Implement all four core modules of the SimCLR architecture in PyTorch without relying on high-level self-supervised frameworks.
    \item \textit{Standard Benchmark Reproduction:} Evaluate SimCLR on the CIFAR-10 benchmark \cite{krizhevsky2009learning} using the frozen linear evaluation protocol, comparing empirical results directly with the authors' published baselines (Appendix B.9 and Figure B.7).
    \item \textit{Controlled Scaling Experiment:} Investigate an open research question untouched by the original paper's ablation studies: how the volume of \textit{unlabeled pretraining data} affects representation quality under constant linear evaluation supervision.
    \item \textit{Interactive Production Deployment:} Deploy the optimized model as an accessible web application featuring Stack Overflow-inspired UI/UX, integrated observability, and sub-20ms inference latency.
\end{enumerate}

\section{Literature Review / Background}
\subsection{The Core Problem and Historical Importance}
The ambition to learn invariant, semantically grounded visual representations without human labels dates back to foundational work on energy-based models and dimensionality reduction \cite{hadsell2006dimensionality}. In classical self-supervised learning, researchers designed artificial ``pretext tasks'' where pseudo-labels could be derived automatically from input structure. Prominent examples include context prediction (predicting the spatial arrangement of adjacent image patches) \cite{doersch2015unsupervised}, solving jigsaw puzzles \cite{noroozi2016unsupervised}, colorizing grayscale images \cite{zhang2016colorful}, and predicting random image rotations \cite{gidaris2018unsupervised}. While effective at forcing the network to learn low-level edges and part relationships, handcrafted pretext tasks suffer from severe limitations: they rely on task-specific heuristics, induce semantic shortcuts, and produce feature representations that fail to generalize broadly across diverse downstream vision tasks.

\subsection{Evolution of Contrastive Learning Approaches}
Contrastive learning departed from discrete pretext classification by formulating representation learning as metric learning in continuous latent space. The fundamental principle is to map inputs to an embedding space such that semantically similar (``positive'') pairs are pulled close together, while dissimilar (``negative'') pairs are pushed far apart.

Before SimCLR, several landmark frameworks explored this formulation:
\begin{itemize}
    \item \textbf{Instance Discrimination \& Memory Banks:} Wu et al. \cite{wu2018unsupervised} proposed treating every image instance in the training set as a distinct class. To store negative representations across optimization steps without computing full-dataset gradients, they introduced an explicit \textit{memory bank} that maintained moving-average representations of all dataset instances. However, the memory bank introduced substantial memory overhead and representational lag, as embeddings stored in the bank became stale as the network weights evolved.
    \item \textbf{Contrastive Predictive Coding (CPC):} van den Oord et al. \cite{oord2018representation} introduced the InfoNCE loss, training autoregressive encoders to predict representations of future or neighboring spatial patches within an image while contrasting against negative patches drawn from other images.
    \item \textbf{Augmented Multiscale Dim (AMDIM):} Bachman et al. \cite{bachman2019learning} maximized mutual information across multiple intermediate feature maps. Although achieving competitive performance, AMDIM required highly complex, specialized multi-scale network architectures.
    \item \textbf{Momentum Contrast (MoCo):} He et al. \cite{he2020momentum} addressed the staleness of memory banks by introducing a dynamic queue decoupled from the dictionary encoder, updated via a momentum-based moving-average parameter tracker.
\end{itemize}

Table \ref{tab:pre_simclr} summarizes the core paradigms preceding SimCLR and highlights their key architectural constraints.

\begin{table}[htbp]
\caption{Approaches Preceding SimCLR and Their Limitations}
\label{tab:pre_simclr}
\centering
\small
\begin{tabular}{lll}
\toprule
\textbf{Approach} & \textbf{Key Idea} & \textbf{Limitation Relevant to SimCLR} \\ \midrule
Pretext Tasks & Predict patch layout / rotation & Relies on designer heuristics \\
Instance Discr. & Treat each image as a class & Requires large, stale memory bank \\
CPC & Predict patches via InfoNCE & Rigid patch splitting and context net \\
AMDIM & Maximize local-to-global MI & Complex, constrained architectures \\
MoCo & Decouple dict. with momentum & Requires auxiliary momentum encoder \\ \bottomrule
\end{tabular}
\end{table}

\subsection{SimCLR: Novelty and Key Contributions}
SimCLR \cite{chen2020simple} simplified contrastive learning by eliminating auxiliary memory banks, specialized pretext networks, and momentum encoders. The authors demonstrated that competitive representations could be learned entirely within standard mini-batches via four unified design choices:
\begin{enumerate}
    \item \textbf{Composition of Data Augmentations:} The authors proved that no single data augmentation strategy suffices. In particular, composing random resized cropping with strong color jittering is critical; without color distortion, the network easily exploits shortcut learning by matching low-level color histograms across cropped views rather than abstract visual semantics.
    \item \textbf{Nonlinear Projection Head:} Introducing a 2-layer MLP projection head $g(\cdot)$ between the base representation $h$ and the contrastive loss function $\mathcal{L}_{\text{NT-Xent}}$ yielded substantial performance gains (+10\% on ImageNet). Crucially, the authors demonstrated that representations before the projection head ($h$) are substantially richer and more linearly separable than representations after the head ($z$), because $z$ discards information invariant to the contrastive objective.
    \item \textbf{Normalized Temperature-Scaled Cross-Entropy (NT-Xent):} A symmetric, normalized cross-entropy formulation incorporating cosine similarity and a tunable temperature parameter $\tau$.
    \item \textbf{Large Mini-Batch Scaling:} By treating all other $2N-2$ augmented views in a batch of $N$ images as negatives, SimCLR directly leverages large mini-batch optimization to provide sufficient negative sampling.
\end{enumerate}

\subsection{Theoretical Foundations: Alignment and Uniformity}
Wang and Isola \cite{wang2020understanding} formalized the mechanics of contrastive loss on the unit hypersphere $\mathcal{S}^{d-1}$, proving that the asymptotic limit of the NT-Xent objective optimizes two distinct geometric properties:
\begin{equation}
\mathcal{L}_{\text{align}} \triangleq \mathbb{E}_{(x, x^{+}) \sim p_{\text{pos}}} \left[ \|f(x) - f(x^{+})\|_{2}^{2} \right]
\end{equation}
\begin{equation}
\mathcal{L}_{\text{uniform}} \triangleq \log \mathbb{E}_{x, y \stackrel{i.i.d.}{\sim} p_{\text{data}}} \left[ e^{-2\|f(x) - f(y)\|_{2}^{2}} \right]
\end{equation}
The temperature parameter $\tau$ serves as a hardness controller: smaller $\tau$ penalizes hard negative instances exponentially more, driving tighter uniformity at the cost of potential representational collapse if negative samples are noisy.

\subsection{Literature on Pretraining Data Scale}
While Chen et al. \cite{chen2020simple} investigated downstream label fraction scaling (evaluating linear probe accuracy with 1\% vs. 10\% vs. 100\% labeled examples on ImageNet), the original paper held the volume of unlabeled pretraining data strictly constant. Subsequent empirical investigations by Goyal et al. \cite{goyal2019scaling} and Cole et al. \cite{cole2022does} highlighted that self-supervised representations exhibit power-law scaling properties with respect to pretraining data volume, but performance gains saturate rapidly unless encoder parameter capacity is concurrently expanded.

\section{Paper Methodology}
The SimCLR framework consists of four sequential modules executed during self-supervised pretraining, followed by a linear evaluation protocol.

\subsection{Stochastic Data Augmentation ($T$)}
Given an unlabeled image $x$, two stochastic data augmentation operators $t, t' \sim \mathcal{T}$ are sampled independently and applied to produce a correlated positive pair of views $(\tilde{x}_i, \tilde{x}_j)$. For CIFAR-10 ($32 \times 32$), the augmentation sequence comprises:
\begin{enumerate}
    \item \textbf{Random Resized Crop:} Area scale $s \in [0.08, 1.0]$, aspect ratio in $[3/4, 4/3]$, bicubic interpolation to $32 \times 32$.
    \item \textbf{Random Horizontal Flip:} Probability $p = 0.5$.
    \item \textbf{Color Jitter:} Probability $p = 0.8$, distortion strength $s=0.5$ (brightness 0.4, contrast 0.4, saturation 0.4, hue 0.1).
    \item \textbf{Random Grayscale:} Probability $p = 0.2$.
    \item \textbf{No Gaussian Blur:} Following the official CIFAR-10 specification (Chen et al. commit \texttt{383d414}), blur is omitted on $32 \times 32$ pixels.
\end{enumerate}

\subsection{Base Encoder ($f(\cdot)$)}
The base visual encoder extracts feature representations $h_i = f(\tilde{x}_i) \in \mathbb{R}^{d}$. We utilize a modified \textbf{ResNet-18} adapted specifically for low-resolution CIFAR-10 inputs:
\begin{itemize}
    \item \textbf{Modified CIFAR Stem:} The standard ImageNet stem ($7 \times 7$ conv, stride 2, plus $3 \times 3$ max-pool) destroys fine spatial details on $32 \times 32$ inputs. We replace this stem with a single $3 \times 3$ convolution (stride 1, 64 channels), batch normalization, and ReLU activation, omitting max-pooling.
    \item \textbf{Residual Stages:} Four basic residual stages ($[2, 2, 2, 2]$ blocks) with feature dimensions $64 \to 128 \to 256 \to 512$.
    \item \textbf{Global Average Pooling:} Collapses spatial dimensions to yield representation vector $h_i \in \mathbb{R}^{512}$.
\end{itemize}

\subsection{Projection Head ($g(\cdot)$)}
A 2-layer multi-layer perceptron (MLP) maps representation $h_i$ into the lower-dimensional latent space:
\begin{equation}
z_i = g(h_i) = W^{(2)} \sigma\left(\text{BN}\left(W^{(1)} h_i\right)\right)
\end{equation}
where $W^{(1)} \in \mathbb{R}^{512 \times 512}$, $\sigma$ denotes ReLU activation, $\text{BN}$ is Batch Normalization, and $W^{(2)} \in \mathbb{R}^{128 \times 512}$ produces latent projection $z_i \in \mathbb{R}^{128}$. The projection head is utilized solely during pretraining and is discarded prior to linear evaluation and inference deployment.

\subsection{NT-Xent Loss Function}
For a mini-batch of $N$ input images, the augmentation pipeline generates $2N$ views. For each anchor view $i$, there exists one positive view $j$, while the remaining $2N-2$ views are treated as negatives. Let $\text{sim}(u, v) = \frac{u^\top v}{\|u\|_2 \|v\|_2}$ denote cosine similarity. The loss for pair $(i, j)$ is:
\begin{equation}
\ell(i, j) = -\log \frac{\exp\left(\frac{\text{sim}(z_i, z_j)}{\tau}\right)}{\sum_{k=1}^{2N} \mathbb{I}_{[k \neq i]} \exp\left(\frac{\text{sim}(z_i, z_k)}{\tau}\right)}
\end{equation}
The total loss is averaged across all positive pairs in the batch:
\begin{equation}
\mathcal{L}_{\text{NT-Xent}} = \frac{1}{2N} \sum_{k=1}^{N} \left[ \ell(2k-1, 2k) + \ell(2k, 2k-1) \right]
\end{equation}

\subsection{Linear Evaluation Protocol}
All parameters of encoder $f(\cdot)$ are frozen (\texttt{requires\_grad = False} and \texttt{model.eval()}). A single linear layer ($W_{\text{cls}} \in \mathbb{R}^{10 \times 512}, b_{\text{cls}} \in \mathbb{R}^{10}$) is trained on representations $h$ using cross-entropy loss under crop and flip augmentations. Performance is reported as Top-1 test accuracy on the 10,000-image test set.

\section{Implementation Details}
\subsection{Software Environment and Dependencies}
The implementation was written in Python 3.9 utilizing native PyTorch (v2.0+) and Torchvision:
\begin{itemize}
    \item \textbf{Core Framework:} \texttt{torch>=2.0.0}, \texttt{torchvision>=0.15.0}
    \item \textbf{Data \& Scientific:} \texttt{numpy>=1.22.0}, \texttt{pandas>=1.5.0}, \texttt{scipy>=1.10.0}
    \item \textbf{Testing:} \texttt{pytest>=7.0.0} (8 automated unit tests)
    \item \textbf{Deployment:} \texttt{gradio==5.20.0}, \texttt{huggingface\_hub<0.26.0}, \texttt{Pillow>=9.5.0}
    \item \textbf{Observability:} \texttt{wandb>=0.15.0}, \texttt{arize-phoenix==12.15.1}, \texttt{graphql-core==3.2.3}
\end{itemize}

\subsection{Compute Hardware and Feasibility Adjustments}
To execute the complete research cycle within an accessible academic budget (an Apple Silicon M4 Pro GPU with 14 cores, 24 GB unified memory, supplemented by Colab T4 GPUs), we instituted four principled feasibility adjustments:
\begin{enumerate}
    \item \textbf{Architecture:} ResNet-18 modified with CIFAR stem (11.17M parameters) instead of ResNet-50 (24.0M parameters).
    \item \textbf{Mini-Batch Size:} $N = 256$ ($2N = 512$ views per step), fitting within local accelerator memory.
    \item \textbf{Pretraining Budget:} 50 epochs (8,750 optimization steps at batch 256 on 45,000 samples).
    \item \textbf{Optimizer:} Replaced LARS with \textbf{Adam} ($lr = 10^{-3}$) after empirical validation on held-out splits.
\end{enumerate}

\subsection{Explicit Code Attribution}
In accordance with Deliverable 3 Guidelines:
\begin{itemize}
    \item \textbf{Team-Implemented (From Scratch):} Augmentation pipeline (\texttt{src/simclr/augment.py}), CIFAR-stem ResNet-18 and MLP projection head (\texttt{src/simclr/model.py}), vectorized NT-Xent loss (\texttt{src/simclr/loss.py}), pretraining and linear evaluation loops (\texttt{src/simclr/train\_*.py}), stratified subset sampler (\texttt{src/simclr/data.py}), unit tests (\texttt{tests/}), and Gradio application (\texttt{app/app.py}).
    \item \textbf{Code Adapted / Referenced:} \texttt{google-research/simclr} (TF2) inspected for CIFAR hyperparameters ($s=0.5, \tau=0.5$); Torchvision \texttt{BasicBlock} adapted into custom \texttt{ResNet18CIFAR}.
    \item \textbf{Results Attribution:} Paper baselines are labeled \texttt{PAPER-REPORTED}; team results are labeled \texttt{MEASURED}.
\end{itemize}

\section{Dataset and Preprocessing / EDA}
\subsection{Dataset Overview}
We utilize the standard \textbf{CIFAR-10} benchmark \cite{krizhevsky2009learning}, consisting of 60,000 $32 \times 32$ color (RGB) images across 10 classes: \textit{airplane, automobile, bird, cat, deer, dog, frog, horse, ship, truck}. Classes are perfectly balanced (5,000 train / 1,000 test per class).

\subsection{Data Quality and Statistical Properties}
Audit checks verified zero missing values, zero corrupted files, and zero duplicate IDs. Analysis of luminance and chromaticity revealed that \textit{airplane} and \textit{ship} have the highest brightness (medians $\approx 142/255$ and $133/255$) and dominant blue channels ($B \approx 150, 141$) due to sky/water backgrounds. Land animals exhibit warm tones ($R > G > B$), and vehicles show neutral color balance.

\subsection{Raw-Pixel PCA Baseline}
Fitting PCA with 50 components on 4,000 raw training images and testing a 5-NN classifier on 1,000 held-out images yielded an accuracy of only \textbf{30.8\%}. This demonstrates that raw-pixel Euclidean distances fail to capture visual semantics, motivating deep representation learning.

\subsection{Data Splitting}
We partitioned the 50,000 training images into a 45,000-image pretraining pool and a 5,000-image stratified validation split (\texttt{data/splits/val\_idx.npy}). The validation split was strictly reserved for pilot tuning; the 10,000-image test set was evaluated only once per final model.

\section{Reproduction Results}
\subsection{Pilot Tuning (Stage S5)}
Five pilot configurations were evaluated on 4,500 images for 5 pretraining + 5 linear evaluation epochs on the validation split (Table \ref{tab:pilots}). Adam ($lr=10^{-3}, \tau=0.5$) outperformed SGD by $+2.70\%$ Top-1 accuracy and was selected for all full runs.

\begin{table}[htbp]
\caption{Pilot Hyperparameter Selection Experiments}
\label{tab:pilots}
\centering
\small
\begin{tabular}{lccccc}
\toprule
\textbf{Pilot} & \textbf{Optimizer} & \textbf{LR} & \textbf{$\tau$} & \textbf{Loss} & \textbf{Val Top-1} \\ \midrule
P1 & SGD & 0.1 & 0.5 & 5.517 & 36.63\% \\
P2 & SGD & 0.3 & 0.5 & 5.497 & 36.67\% \\
\textbf{P3 (Selected)} & \textbf{Adam} & \textbf{1e-3} & \textbf{0.5} & \textbf{5.462} & \textbf{39.33\%} \\
P4 & AdamW & 1e-3 & 0.5 & 5.462 & 39.06\% \\
P5 & SGD & 0.3 & 0.1 & 5.270 & 37.10\% \\ \bottomrule
\end{tabular}
\end{table}

\subsection{Main Reproduction Run (Stage S6)}
Pretraining for 50 epochs on 45,000 images followed by 40 linear evaluation epochs yielded:
\begin{itemize}
    \item \textbf{Test Top-1 Accuracy:} \textbf{75.23\%} (7,523 / 10,000).
    \item \textbf{Test Top-5 Accuracy:} \textbf{98.71\%} (9,871 / 10,000).
    \item \textbf{Test Macro-F1:} \textbf{0.7512}.
    \item \textbf{Pretraining Contrastive Accuracy:} \textbf{62.35\%}.
\end{itemize}

\subsection{Comparison with Published Baselines}
Table \ref{tab:comp_main} presents the side-by-side comparison.

\begin{table}[htbp]
\caption{Comparison Against Published Baselines}
\label{tab:comp_main}
\centering
\small
\begin{tabular}{lccc}
\toprule
\textbf{Parameter / Metric} & \textbf{Chen et al. App. B.9} & \textbf{Official Repo} & \textbf{Reproduction} \\ \midrule
Backbone & ResNet-50 & ResNet-18 & ResNet-18 \\
Parameters & 24.0M & 11.2M & 11.2M \\
Epochs & 1,000 & 1,000 & \textbf{50} \\
Batch Size & 1,024 & 512 & \textbf{256} \\
Optimizer & LARS & LARS & \textbf{Adam} \\
Total Steps & $\sim$48,800 & $\sim$97,600 & \textbf{8,750} \\ \midrule
\textbf{Test Top-1} & \textbf{94.0\%} & \textbf{$\approx$91.0\%} & \textbf{75.23\%} \\
100-Epoch Est. & $\approx$84.0\% (Fig. B.7) & $\approx$83.5\% & \textbf{75.23\% (50 ep)} \\ \bottomrule
\end{tabular}
\end{table}

\section{Additional Experiment: Pretraining Data Scaling}
\subsection{Research Question}
How does the volume of unlabeled pretraining data affect representation quality when downstream labeled supervision is held constant?

\subsection{Methodology}
We extracted nested stratified subsets: 10\% (4,500), 25\% (11,250), 50\% (22,500), and 100\% (45,000). Each subset was pretrained for 50 epochs under Adam ($lr=10^{-3}, \tau=0.5$). Each frozen checkpoint was evaluated on the \textbf{exact same full labeled dataset} (45,000 images) for 40 linear evaluation epochs.

\subsection{Results}
Table \ref{tab:scaling} demonstrates monotonic scaling with diminishing marginal returns.

\begin{table}[htbp]
\caption{Pretraining Data Fraction vs. Downstream Accuracy}
\label{tab:scaling}
\centering
\begin{tabular}{lcccc}
\toprule
\textbf{Fraction} & \textbf{Samples} & \textbf{Top-1 Acc} & \textbf{Top-5 Acc} & \textbf{Marginal Gain} \\ \midrule
10\% & 4,500 & 54.17\% & 94.52\% & Baseline \\
25\% & 11,250 & 65.37\% & 97.14\% & +11.20\% \\
50\% & 22,500 & 71.04\% & 98.07\% & +5.67\% \\
100\% & 45,000 & 75.23\% & 98.71\% & +4.19\% \\ \bottomrule
\end{tabular}
\end{table}

\section{Analysis and Discussion}
\subsection{Reproduction Gap Attribution}
The gap between our 75.23\% result and the paper's 94.0\% is fully explained by compute and architectural differences:
\begin{enumerate}
    \item \textbf{Optimization Steps:} 50 epochs (8,750 steps) vs. 1,000 epochs (48,800 steps). Chen et al. Figure B.7 shows that at 100 epochs, ResNet-50 achieves $\approx 84\%$; reaching 75.23\% at 50 epochs with a smaller ResNet-18 aligns closely with expected convergence curves.
    \item \textbf{Model Capacity:} ResNet-18 (11.2M params) vs. ResNet-50 (24.0M params).
    \item \textbf{Negative Sampling:} Batch size 256 provides 510 negative views per step versus 2,046 at batch size 1024, slightly softening uniformity pressure on the hypersphere.
\end{enumerate}

\subsection{Power-Law Scaling}
The data scaling ablation exhibits clear logarithmic behavior: transitioning from 10\% to 25\% data yields $+11.20\%$, whereas quadrupling data from 25\% to 100\% yields $+9.86\%$ cumulative gain. This confirms theoretical predictions that unlabeled data volume provides high marginal value initially, but saturates without expanding model parameter capacity.

\section{Model Deployment}
\subsection{Interactive Web Application}
We developed and deployed an interactive web application built with Gradio and styled after Stack Overflow's design aesthetic:
\begin{itemize}
    \item \textbf{Hosting:} Hugging Face Spaces (\url{https://huggingface.co/spaces/MianGenz/Sim_clr}).
    \item \textbf{Direct URL:} \url{https://miangenz-sim-clr.hf.space}.
    \item \textbf{Artifact:} \texttt{app/simclr\_cifar10\_eval.pt} (42.7 MB, SHA-256: \texttt{12358c35...}).
    \item \textbf{Input/Output:} Accepts standard images, resizes to $32 \times 32$, normalizes with CIFAR statistics, runs frozen encoder + linear probe, and outputs Top-3 probabilities and inference latency ($18.15$ ms average).
\end{itemize}

\subsection{Observability}
Request tracing was instrumented via Arize Phoenix (OpenTelemetry), and training metrics were logged via Weights \& Biases.

\section{Limitations}
\begin{itemize}
    \item Pretraining was bounded at 50 epochs due to single-accelerator constraints.
    \item Mini-batch size was restricted to 256 without multi-GPU synchronization.
    \item CIFAR stem is specialized for $32 \times 32$ inputs; high-resolution images must be downsampled.
\end{itemize}

\section{Conclusion}
We successfully reproduced SimCLR (v1) from scratch in PyTorch on CIFAR-10, attaining 75.23\% Top-1 test accuracy in 50 epochs on a single GPU. Our scaling experiment proved that unlabeled pretraining data follows logarithmic scaling with diminishing marginal returns. Finally, the model was successfully deployed with Stack Overflow-inspired UI and production observability on Hugging Face Spaces.

\begin{thebibliography}{00}
\bibitem{chen2020simple} T. Chen, S. Kornblith, M. Norouzi, and G. Hinton, ``A simple framework for contrastive learning of visual representations,'' in \textit{Proc. 37th Int. Conf. Machine Learning (ICML)}, PMLR 119, 2020, pp. 1597--1607.
\bibitem{krizhevsky2009learning} A. Krizhevsky and G. Hinton, ``Learning multiple layers of features from tiny images,'' Univ. Toronto, Tech. Rep., 2009.
\bibitem{russakovsky2015imagenet} O. Russakovsky et al., ``ImageNet large scale visual recognition challenge,'' \textit{Int. J. Computer Vision}, vol. 115, no. 3, pp. 211--252, 2015.
\bibitem{he2016deep} K. He, X. Zhang, S. Ren, and J. Sun, ``Deep residual learning for image recognition,'' in \textit{Proc. IEEE Conf. Computer Vision and Pattern Recognition (CVPR)}, 2016, pp. 770--778.
\bibitem{hadsell2006dimensionality} R. Hadsell, S. Chopra, and Y. LeCun, ``Dimensionality reduction by learning an invariant mapping,'' in \textit{Proc. CVPR}, vol. 2, 2006, pp. 1735--1742.
\bibitem{doersch2015unsupervised} C. Doersch, A. Gupta, and A. A. Efros, ``Unsupervised visual representation learning by context prediction,'' in \textit{Proc. IEEE Int. Conf. Computer Vision (ICCV)}, 2015, pp. 1422--1430.
\bibitem{noroozi2016unsupervised} M. Noroozi and P. Favaro, ``Unsupervised learning of visual representations by solving jigsaw puzzles,'' in \textit{Proc. ECCV}, 2016, pp. 69--84.
\bibitem{zhang2016colorful} R. Zhang, P. Isola, and A. A. Efros, ``Colorful image colorization,'' in \textit{Proc. ECCV}, 2016, pp. 649--666.
\bibitem{gidaris2018unsupervised} S. Gidaris, P. Singh, and N. Komodakis, ``Unsupervised representation learning by predicting image rotations,'' in \textit{Int. Conf. Learning Representations (ICLR)}, 2018.
\bibitem{wu2018unsupervised} Z. Wu, Y. Xiong, S. X. Yu, and D. Lin, ``Unsupervised feature learning via non-parametric instance discrimination,'' in \textit{Proc. CVPR}, 2018, pp. 3733--3742.
\bibitem{oord2018representation} A. van den Oord, Y. Li, and O. Vinyals, ``Representation learning with contrastive predictive coding,'' \textit{arXiv preprint arXiv:1807.03748}, 2018.
\bibitem{bachman2019learning} P. Bachman, R. D. Hjelm, and W. Buchwalter, ``Learning representations by maximizing mutual information across views,'' in \textit{NeurIPS}, 2019, pp. 15509--15519.
\bibitem{he2020momentum} K. He, H. Fan, Y. Wu, S. Xie, and R. Girshick, ``Momentum contrast for unsupervised visual representation learning,'' in \textit{Proc. CVPR}, 2020, pp. 9729--9738.
\bibitem{wang2020understanding} T. Wang and P. Isola, ``Understanding contrastive representation learning through alignment and uniformity on the hypersphere,'' in \textit{Proc. ICML}, PMLR 119, 2020, pp. 9929--9939.
\bibitem{goyal2019scaling} P. Goyal, D. Mahajan, A. Gupta, and I. Misra, ``Scaling and benchmarking self-supervised visual representation learning,'' in \textit{Proc. ICCV}, 2019, pp. 6391--6400.
\bibitem{cole2022does} E. Cole, X. Yang, K. Wilber, O. Mac Aodha, and S. Belongie, ``When does contrastive visual representation learning work?,'' in \textit{Proc. CVPR}, 2022, pp. 14755--14764.
\end{thebibliography}

\end{document}
